\paragraph{From the Herr pairs to normalized colength.}
\label{par:row-colength}
On each component use the module \(I_B\) of exceptional coordinates
from Proposition~\ref{prop:row-exceptional-coordinates}. Choose lifts
of an \(O\)-basis of \(I_B\).
Apply \(\alpha\) to each lift, choosing a representative invariant
modulo \(XN_c\). By the order-raising argument in
Section~\ref{par:row-metric}, its \(p^\sigma\)-multiple can be
corrected modulo \(XN_c\) to an invariant to order \(M\).
Complete it by the unique deep second Herr coordinate. This gives
at most \(e(e+1)m\) fixed finite classes on this component.
Write \(H^1_f(K,T)\) for the inverse image of the rational
finite condition in integral cohomology, including torsion.
Write \(J\) for the module generated by all these classes on all
components, in the quotient
\[
              C_F=H^1_f(K,T_2)/F_*H^1_f(K,T_1).
\]
There are at most \(ne(e+1)m=D(e+1)m\) generators of \(J\).

\begin{proposition}[Normalized colength of the finite local condition]\label{prop:local-row-colength}
Let \(C_F\) be the cokernel on integral finite cohomology
of the map of Definition~\ref{hyp:first-level-row}. With
\(J\) the span of the exceptional classes constructed from
Proposition~\ref{prop:row-exceptional-coordinates}, the
normalized length satisfies
\[
 \frac{\operatorname{length}_O C_F}{\operatorname{length}_O(O/pO)}
       \le 2(D+1)rY+D(e+1)m\,18h .
\]
For \(p\ge3\), \(h\ge1\) and \(e\in\{1,2\}\), the common exponent
\(Y\) below is at most \(80(eh/p+eS+e\log_p h+1+\eta)\).
The additional \(18h\) allowance occurs on at most \(D(e+1)m\)
exceptional directions.
\end{proposition}

\begin{proof}
\proofstep{Lifting modulo the fixed exceptional classes.}
Let \((w,z)\) be a target Herr pair with \(w\in E^{-H}P_c\)
invariant to order \(M\) and \(z\in X^MN\). In the model
\(\beta(w)\), apply \(P(\gamma)\) and use the splitting into
main and residual coordinates. Expand the residual modulo
\(XN^{\rm mod}\) in the chosen basis of \(I_B\).
Return from the model to the Wach lattice using \(\alpha\). The integral multiplier and the cost of the two comparison maps yield
\[
 \frac{p^\kappa}{P(1)}\,P(1)p^{2\sigma}w=p^{2\sigma+\kappa}w.
\]

At exponent \(y_1\) every main tensor on its right lifts from
the source modulo \(XN_c\), by Section~\ref{par:row-mixed}.
The residual terms are integral combinations of the fixed
representatives. Multiply by \(p^\sigma\), raise all these
invariants to precision \(M\) and complete the source pair.
The difference of the first coordinates is invariant to order \(M\)
and lies in \(XN_c\). It lies in \(X^MN_c\) by uniqueness of
invariant lifting. Naturality and the deep coboundary criterion now give
\[
              p^{\,3\sigma+\kappa+y_1}[(w,z)]\in J
                       \quad\hbox{in } C_F .
\]
In particular the precision raising costs one \(\sigma\) for the
entire relation, irrespective of its number of terms.

\proofstep{Reduction of a finite class to the tested form.}
It remains to put a general finite target class into this form.
Its extension of the unit by \(T_2\) is crystalline over the fraction
field. Lift \(p^\rho\) times the unit generator in each component
to \(x_c\) in the extension Wach lattice. Its cyclotomic difference
lies in \(XN_c(T_2)\) since the quotient generator is fixed.
Here one needs the kernel intersection: the integral étale lattice
sequence is exact and the subrepresentation has the induced Wach
lattice. The same intersection holds in the Frobenius span, as the
coefficient ring is free over its Frobenius image. In the bounded
pole module it follows after multiplying by \(E^H\).
The graded correction produces \(x'_c\) with
\[
 x'_c-p^\sigma x_c\in XN_c(T_2),\qquad (\gamma-1)x'_c\in X^M N_c(T_2).
\]

The correction uses the graded scalar calculation on that kernel.
Its projection is still \(p^{\rho+\sigma}\) times the unit.
Consequently
\[
                        w_c=\varphi x'_{c-1}-x'_c
\]
belongs to the kernel in \(E^{-H}P_c\) and is invariant to order
\(M\), since \(\gamma\) commutes with \(\varphi\) and \(\varphi\)
preserves \(X^MN\). Choose the unique \(z_c\in X^MN\) with
\((\varphi-1)z_c=(\gamma-1)w_c\). On summing in \(c\), uniqueness
gives \(\sum z_c=(\gamma-1)\sum x'_c\), since both sides are deep
and have the same Frobenius difference. Thus
\(\sum(w_c,z_c)\) represents \(p^{\rho+\sigma}\) times the original
class by the connecting map, allowing its fixed unit. We have proved
\[
       p^Y C_F\subset J,\qquad
       Y=2\rho+4\sigma+\kappa+\max(2\sigma,y_c).
       \tag{R0}\label{eq:tag-r0}
\]

\proofstep{Bounding the exceptional submodule.}
A uniform exponent on \(J\) follows from the coarser full-level lift.
Over \(K_0(\mu_p)\) the two height bounds are \(h,h\), and the jumps
of sources and kernels remain in \([-h,h]\). Restriction and
corestriction cost the degree \(s\), a unit in \(O\).
In the full-level computation of Section~\ref{subsec:ramified-lift}
the two Wach losses are at most \(4h+v_p((4h)!)\) each.

The
kernel correction has loss at most
\[
 h+v_p(h!)+\sum_{j\in\mathcal J}\bigl(1+v_p(j)\bigr),
 \qquad
 \mathcal J\subset[-h,2h]\cap\mathbb Z,\quad 0\notin\mathcal J .
\]
Indeed the shift of the jump set uses degrees \(0,\ldots,h\) and
the positive degrees for the trivial component are at most \(h\);
zero eigenvalue differences are discarded in the rational coboundary
test. Even using the full set of nonzero signed integers in this
interval gives valuation sum at most \(v_p(h!)+v_p((2h)!)\).
Since \(p\ge3\), \(v_p(j!)\le j/2\). The total exponent is at most
\(18h\). Thus \(p^{18h}\) kills \(J\).

For clarity about the other rank factor, the rational Bloch--Kato
dimension formula bounds \(\dim H^1_f(K,V_2)\) by
\[
 \dim H^0(K,V_2)+\dim(D_{\rm dR}(V_2)/{\rm Fil}^0) \le (D+1)r.
\]
For a uniformizer \(\varpi\) of \(O\), the
cohomology sequence of multiplication by \(\varpi\) exhibits
\(H^1(K,T_2)[\varpi]\) as a quotient of \(H^0(K,T_2/\varpi T_2)\).
Its residue dimension is at most \(r\). Local integral cohomology
is finite type by local Tate finiteness. Its torsion and the
torsion of its finite-condition submodule therefore need at most
\(r\) generators by the structure theorem over a DVR.
Hence \(2(D+1)r\) generators suffice for the finite-condition
module and its quotients. Taking lengths first in \(C_F/J\) and
then in \(J\) gives
\[
 \frac{\operatorname{length}_O C_F}{\operatorname{length}_O(O/pO)}
     \le 2(D+1)rY+D(e+1)m\,18h .
\]
In particular the loss proportional to \(h\) per generator only
occurs on the \(D(e+1)m\) exceptional directions.

\proofstep{The common scalar exponent.}
For the numerical comparison put \(v=H/s\). Then
\(v\le 3eh/(2p)+eS/2+1\), as \(p\ge3\) and \(e\le p-1\).
The factorial bound gives \(\rho\le6v+3/2\) and \(\sigma\le6v\).
Therefore
\[
 Y\le48v+3+e(3S/2+\log_p H)+\eta+3/2+\log_p H .
\]

For \(h\ge1\), \(h+S+1\le h p^{S+1}\), so this is at most
\(80(eh/p+eS+e\log_p h+1+\eta)\) for \(e=1,2\).
The host always has \(h\ge1\). Substitution gives the ramified
estimate above and the same formula with \(e=1\) on an unramified
one-sign component. All choices precede the finite target class
apart from its own lift and expansion. The constants are independent
of the number of coefficient places, their ramification, the number
of Frobenius components, the source rank, and the outer jumps.
\label{cor:row-colength-1}
\end{proof}
